\section{Introduction}
\label{sec:intro}

The EU AI Act (Regulation 2024/1689) anchors conformity assessment on a system's
\emph{intended purpose} (Art.~9, Annex~IV)~\citep{euAIAct2024}. For deployed AI
agents, the question of whether an agent pursues its intended purpose is not
answered by capability benchmarks or static alignment tests: it requires reasoning
about what the agent is actually optimising for across sequences of real tool
calls. An agent that is told to ``find the cheapest option'' but instead maximises
the number of tool calls it makes is not failing a capability test; it is pursuing
a different objective from the one it declares.

The literature that bears on this question divides into three clusters that do not
overlap. First, agent red-teaming and misalignment benchmarks
\citep{agentMisalignment2025} probe capability and misalignment propensity using
fixed attack catalogues; they establish that deployed agents can and do behave in
misaligned ways, but they do not infer what an agent is optimising for, and they do
not use such an inference to steer probing. Second, IRL for alignment auditing
\citep{alignmentAuditor2025, ir3contrastive, insightsInverse2024} reconstructs objectives
from model weights or RLHF training signals; these methods are not applicable to
deployed black-box agents accessible only through an API. Third, goal inference
from behaviour \citep{sips2020, communicatingAgents2023} operates in structured
planning domains rather than in the unstructured tool-call space of deployed
agents.

\paragraph{The empty cell.}
No existing tool combines (a) IRL-based objective inference from tool-call
trajectories, (b) divergence computation from a declared purpose, and (c)
adversarial probe steering driven by that divergence, whilst emitting the result as
regulatory evidence. We introduce the Inverse-Objective Regulator (IOR) to occupy
this cell.

\paragraph{Contributions.}
\begin{itemize}[leftmargin=1.4em]
  \item IOR, a system design that connects goal inference to adversarial probe
  targeting at the deployed-agent level through a closed loop in which inference
  steers probing and probe results refine inference. We present the loop as a
  design and evaluate its components; closing it empirically is future work
  (Section~\ref{sec:limitations}).
  \item The Goal Decomposition Judge, which bridges natural-language declared
  purposes into the feature space that IRL requires without access to model
  weights, together with a pluggable feature basis spanning a deterministic
  model-free floor, a single judge, and a judge ensemble with quantified
  agreement.
  \item An open gym of three planted-divergence reference agents with deterministic
  seeds, enabling reproducible evaluation of objective recovery.
  \item Divergence reports mapped to OWASP ASI, MITRE ATLAS, and NIST AI RMF nodes,
  making findings directly usable in compliance workflows.
\end{itemize}

\paragraph{Scope.}
IOR operates black-box (API-level traces only, no weights or logits). It audits and
certifies; it does not enforce inline. Because IRL is ill-posed, IOR emits a
distribution over candidate objectives with uncertainty, never a single
overconfident point estimate.
